The final transfer evaluation takes all three ft09-trained seeds of each
algorithm's baseline, novelty and final memory stage (DQN no-change memory;
PPO memory + novelty) and runs them unchanged on
cd82, tr87 and vc33. Each checkpoint plays five rollouts per game, with sampling
seeds 100--104 and the same 1,000-action budget: 270 episodes across six stages.
Five rollouts rather than 20 keep the evaluation cost manageable while retaining
all three training seeds. No held-out outcome influenced model selection or
parameters. This tests transfer from a single training game, not online
adaptation or multi-game training.

All PPO stages and both improved DQN stages complete zero levels on all three
held-out games. DQN baseline seed 1 completes one tr87 level in one of five
rollouts (one of 15 pooled baseline rollouts on tr87); its other seeds complete
none. No agent wins a game, and every episode uses its full action budget. The
single completion does not establish generalisation. Development success on one
game therefore gives little assurance of transfer to a new puzzle under this
protocol. Per-seed results and the frozen protocol are in the evidence folder.
